The annotation of rare cell types in single-cell RNA-sequencing (scRNA-seq) data is among the most consequential and methodologically challenging problems in single-cell biology. Rare populations---spanning plasmacytoid dendritic cell precursors, non-classical monocytes, innate lymphoid cells, and tissue-resident endothelial subtypes---often drive critical immunological and homeostatic processes while representing fewer than 1--5\% of cells in typical tissue atlases \citep{villani2017single,tabulasapiens2022}. Misidentifying or entirely missing these populations distorts downstream differential expression analyses, trajectory inferences, and cell-communication predictions, with cascading effects on biological conclusions \citep{luecken2021benchmarking}.

Semi-supervised deep generative models, exemplified by scANVI \citep{xu2021probabilistic}, represent a powerful approach to cell-type annotation. scANVI learns a shared low-dimensional latent space across batches and donors, leverages partial cell-type labels to shape this space, and applies an inductive classifier to annotate query cells from unseen batches. This inductive setting---where the model is trained on reference cells and applied to held-out query cells without retraining---is the operationally relevant scenario for most atlas-mapping applications.

\paragraph{The rare-cell miscalibration problem.}
Despite its strengths, scANVI's softmax classifier inherits the label imbalance of the training set. When a rare cell type contributes only 5--20 labeled examples while majority types contribute hundreds or thousands, the classifier's posterior probability for the rare class is systematically suppressed---a well-documented consequence of optimizing cross-entropy loss on imbalanced data \citep{xu2023scbalance}. This failure is not captured by overall accuracy metrics, which remain high even when rare-class F1 approaches zero. Critically, the problem is amplified under batch-held-out evaluation: rare cells in the query batch were absent from any training batch, so the model has no within-batch evidence to counter the imbalanced prior.

What makes this situation remediable is a geometric observation: even when the scANVI softmax for the rare class collapses, the latent embedding of rare cells often retains its distinctiveness. The rare prototype remains well-separated from majority prototypes in the learned latent space, even as the classifier ignores this geometry. This disconnect---reliable embeddings but unreliable probabilities---motivates a geometric approach to post-hoc rescue.

\paragraph{Gap in existing methods.}
Prior work on rare-cell identification takes one of three main directions. \emph{De novo} discovery methods \citep{jiang2016giniclust,jindal2018fire,fang2021gapclust,xu2024scCAD} identify small clusters without requiring labeled rare examples, but cannot leverage the few labels that \emph{are} available or produce reliable annotations for known cell types. \emph{Reference-based annotation} methods \citep{aran2019reference,kiselev2018scmap,dominguez2022cross,xu2021probabilistic,lotfollahi2022mapping} optimize aggregate accuracy and scale well to large atlases, but do not address the specific failure mode of rare-class miscalibration. \emph{Imbalance-aware retraining} methods \citep{mulvey2021scsyno,xu2023scbalance} improve minority-class recall by augmenting the training distribution, but require modifying the training pipeline and cannot be applied post-hoc to existing scANVI models.

No existing approach combines geometry-guided rescue---exploiting the fact that rare cells may be geometrically separable even when the classifier is miscalibrated---with marker-gene verification in a fully post-hoc, inductive framework that also provides a pre-deployment diagnostic for when rescue is expected to help.

\paragraph{scRareRefine.}
We present \textbf{scRareRefine} (\cref{fig:framework}), a three-component pipeline that operates downstream of any scANVI-based annotation without retraining. Given the scANVI latent embedding and softmax output for a query dataset, scRareRefine:
\begin{enumerate}[leftmargin=*,label=\textbf{\arabic*.}]
    \item Computes the \emph{separability ratio} $S = \delta_r / \rho_r$ from labeled training cells, where $\delta_r$ is the distance from the rare-class prototype to the nearest majority-class prototype and $\rho_r$ is the intra-class radius of the rare training cluster. This scalar, computable before any query inference, serves as a pre-deployment diagnostic: $S > 1.3$ indicates that the rare cluster is well-separated from majority classes, making prototype-based rescue geometrically valid.

    \item Flags \emph{rescue candidates} by ranking each query cell by its Euclidean distance to all class prototypes. Cells whose nearest-class prototype is the rare class (or within a margin of the top prediction) and whose prediction does not already assign the rare class are selected as candidates.

    \item Confirms rescue via a \emph{marker-gene signature} derived entirely from labeled training data. For each candidate, a score measuring expression of the top discriminative genes for the rare class is compared against the expression of the predicted-class genes. Only candidates passing a validation-tuned threshold are assigned the rare-class label.
\end{enumerate}

When the separability ratio $S < 1.1$, the rare class overlaps with majority classes in latent space, and the candidate set is emptied before marker verification---no rescue occurs and the scANVI prediction is preserved. This structural fallback prevents over-rescue in settings where prototype geometry provides no reliable signal.

\paragraph{Contributions.}
\begin{itemize}[leftmargin=*]
    \item The \textbf{separability ratio}, a scalar diagnostic computable from labeled training data before query inference, that reliably predicts whether prototype-based rescue will be effective.
    \item The \textbf{scRareRefine pipeline}, integrating prototype scoring, a geometric gate, and marker-gene verification under a strict inductive evaluation protocol that prevents any leakage of test-set labels into the rescue procedure.
    \item Empirical evidence across seven datasets showing consistent gains of $+23$--$+78$ pp in rare-class F1 for cell types with $S > 1.3$ and scANVI baseline F1 below 0.75, and near-zero modification (mean $|\Delta\text{F1}| < 0.02$) for cell types outside this regime.
    \item A \textbf{data-efficiency analysis} demonstrating that the rescue advantage is greatest at annotation budgets of five or fewer labeled rare cells---the regime most relevant to discovery settings---with cDC1 F1 improving from 0.003 to 0.986 at $n_\text{rare} = 5$.
\end{itemize}

See \cref{sec:appendix} for implementation details and additional ablation results.
